% Irodov, user-supplied PDF, original printed pp.45--46, PDF pp.54--55.
% Rigid body: rotating bearings, gyroscopic precession, and rail loading.

\exproblem{281}{45}
一根质量 $m=1.40\,\mathrm{kg}$、长度
$\ell_0=1.00\,\mathrm m$ 的均匀细杆水平放置，
一端 $A$ 固定在竖直转轴 $OO'$ 的中点，
可随竖轴以角速度 $\omega$ 匀速转动。
轴上下两个轴承距离为
$\ell=55\,\mathrm{cm}$。
求当下轴承所受水平力为零时所需角速度，
并求此时上轴承所受水平力大小。
\exsolution
杆的质心距轴 $\ell_0/2$，需要水平向心力
$F_c=m\omega^2\ell_0/2$。
杆重力关于轴中点产生大小
$N_g=mg\ell_0/2$ 的倾覆力矩；
上下轴承的水平反力形成力偶以平衡它。
当下轴承的水平反力为零时，
上轴承的水平力 $F_u$ 必须同时满足
\[
F_u=F_c=\frac12m\omega^2\ell_0,
\qquad
F_u\frac{\ell}{2}=mg\frac{\ell_0}{2}.
\]
联立得到
\[
\boxed{\omega=\sqrt{\frac{2g}{\ell}}
\simeq5.97\,\mathrm{rad/s}\approx6.0\,\mathrm{rad/s}},
\]
\[
\boxed{F_u=\frac{mg\ell_0}{\ell}
\simeq25.0\,\mathrm N}.
\]
\excheck 杆固定在竖轴的\emph{中点}，
上轴承到固定点的力臂是 $\ell/2$，
不是全部轴承间距 $\ell$；
采用正确力矩臂后得到 $25\,\mathrm N$，
与底本书末答案一致。

\exproblem{282}{45}
长为 $\ell$、质量为 $m$ 的均匀细杆中点 $C$
固定于竖直轴 $OO'$，
杆与轴夹角为 $\theta$（原图1.71）。
系统绕竖轴以恒角速度 $\omega$ 旋转。
求：(a) 杆关于中点的角动量矢量大小、方向及其竖轴分量；
(b) 半周内角动量矢量增量的模；
(c) 轴承必须提供的外力矩大小。
\exsolution
杆对质心的转动惯量张量沿杆轴分量近似为零，
垂直杆轴的分量
$I_\perp=m\ell^2/12$。
记杆方向单位矢量为 $\bm e$，
竖直单位矢量为 $\bm k$，
转动角速度矢量 $\bm\omega=\omega\bm k$。
关于中点的角动量
\[
\bm L_C=I_\perp[
\bm\omega-(\bm\omega\cdot\bm e)\bm e]
=I_\perp\omega(\bm k-\cos\theta\,\bm e).
\]
其方向在杆与竖轴张成的平面内、垂直于杆，
大小和竖轴投影为
\[
\boxed{|\bm L_C|=\frac{m\ell^2\omega}{12}\sin\theta},
\qquad
\boxed{L_z=\frac{m\ell^2\omega}{12}\sin^2\theta}.
\]
随系统转过半周，角动量的水平分量反向，
竖直分量不变。水平分量大小为
$I_\perp\omega\sin\theta\cos\theta$，
故
\[
\boxed{|\Delta\bm L_C|
=2I_\perp\omega|\sin\theta\cos\theta|
=\frac{m\ell^2\omega}{12}|\sin2\theta|}.
\]
该水平分量以角速度 $\omega$ 旋转，
所以轴承需要的力矩
\[
\boxed{N=\omega|L_{\perp,z}|
=\frac{m\ell^2\omega^2}{24}|\sin2\theta|}.
\]
\excheck 这里 $L_{\perp,z}$ 表示角动量相对于
\emph{竖轴}的水平分量。
(c) 与本书第269题同一几何结构的连续积分应一致；
本题书末给 $\frac1{24}$，而第269题简答给 $\frac12$，
进一步表明后者存在系数印误。

\exproblem{283}{45}
一个陀螺质量为 $m=0.50\,\mathrm{kg}$，
自转轴与竖直夹角 $\theta=30^\circ$，
绕自转轴惯量 $I=2.0\,\mathrm{g\,m^2}$，
自转角速度 $\omega=350\,\mathrm{rad/s}$。
支点到质心的距离 $\ell=10\,\mathrm{cm}$。
在慢进动近似下求：(a) 进动角速度；
(b) 支点对陀螺的水平反力大小和方向。
\exsolution
自转主角动量模为 $L_s=I\omega$。
重力关于支点的力矩大小为
$N_g=mg\ell\sin\theta$。
慢进动时 $\bm L_s$ 以角速度 $\Omega$
绕竖直方向转动，
其矢量变化率大小为 $\Omega L_s\sin\theta$。
因此
\[
\Omega I\omega\sin\theta=mg\ell\sin\theta,
\qquad
\boxed{\Omega=\frac{mg\ell}{I\omega}
\simeq0.701\,\mathrm{rad/s}}.
\]
质心绕竖轴以半径 $\ell\sin\theta$
作匀速圆周运动，水平向心加速度为
$\Omega^2\ell\sin\theta$。
重力无水平分量，故支点水平反力为
\[
\boxed{F_h=m\Omega^2\ell\sin\theta
\simeq1.23\times10^{-2}\,\mathrm N},
\]
方向\emph{水平指向进动圆的旋转轴}。
\excheck 所求是慢进动模型下的支点反力，
不等于陀螺所受重力大小；
原书将其量级约记为 $10\,\mathrm{mN}$。

\exproblem{284}{46}
一陀螺由半径 $R=5.0\,\mathrm{cm}$ 的
均匀圆盘及长 $\ell=10\,\mathrm{cm}$ 的轻杆组成，
杆一端铰接在竖直向上以
$a_0=2.0\,\mathrm{m/s^2}$ 加速的电梯地板上（图1.73）。
圆盘以稳定角速度
$n=0.5$ 转/秒进动。
忽略摩擦及杆质量，求圆盘自身的角速度 $\omega$。
\exsolution
电梯参考系的等效重力加速度为
$g_{\rm eff}=g+a_0$。
均匀圆盘自转惯量
$I=\frac12mR^2$，主角动量
$L_s=I\omega$。
以 $\ell$ 为力臂，重力与惯性力合力矩
$mg_{\rm eff}\ell$ 引起慢进动，
其角速度 $\Omega=2\pi n$，故
\[
\Omega I\omega=mg_{\rm eff}\ell.
\]
消去圆盘质量得
\[
\boxed{\omega=
\frac{2(g+a_0)\ell}{R^2(2\pi n)}
=\frac{(g+a_0)\ell}{\pi nR^2}
\simeq3.01\times10^2\,\mathrm{rad/s}}.
\]
\excheck 电梯\emph{向上}加速使等效重力变大，
因此公式必须含 $g+a_0$；原书简答与此一致。

\exproblem{285}{46}
陀螺质量 $m=1.0\,\mathrm{kg}$，
自转轴转动惯量 $I=4.0\,\mathrm{g\,m^2}$、
自转角速度 $\omega=310\,\mathrm{rad/s}$。
支点固定在一水平加速的物块上，
物块加速度大小为 $a_0=1.0\,\mathrm{m/s^2}$。
支点到质心距离 $\ell=10\,\mathrm{cm}$。
求陀螺稳定慢进动角速度的模及进动轴方向。
\exsolution
在随支点加速的非惯性系中，
除重力 $m\bm g$ 外还有与支点水平加速方向相反的
惯性力 $-m\bm a_0$。
二者合成有效重力场
$\bm g_{\rm eff}=\bm g-\bm a_0$，
其大小
$g_{\rm eff}=\sqrt{g^2+a_0^2}$。
慢进动时
$\Omega I\omega=m\ell g_{\rm eff}$，
所以
\[
\boxed{\Omega=\frac{m\ell\sqrt{g^2+a_0^2}}{I\omega}
\simeq0.796\,\mathrm{rad/s}}.
\]
进动轴与有效重力方向一致，
其相对竖直的倾角为
\[
\boxed{\delta=\arctan(a_0/g)
\simeq5.82^\circ\approx6^\circ},
\]
从竖直方向朝\emph{支点加速度的反方向}倾斜。
\excheck 求进动速率应用有效重力场的\emph{模}，
求进动轴方向则保留其矢量方向。

\exproblem{286}{46}
质量 $m=5.0\,\mathrm{kg}$、
半径 $R=6.0\,\mathrm{cm}$ 的均匀实心球，
绕穿过球心的水平轴以
$\omega=1250\,\mathrm{rad/s}$ 自转。
其水平轴由两个间距 $\ell=15\,\mathrm{cm}$
的轴承支撑。整个支架绕竖直轴以
$\Omega=5.0\,\mathrm{rad/s}$ 匀速转动。
求轴承受到的陀螺力大小及方向关系。
\exsolution
均匀实心球自转惯量为
$I=\frac25mR^2$，主自转角动量大小 $L=I\omega$。
自转轴水平，其角动量方向随支架绕竖直轴以
角速度 $\Omega$ 转动，
因此角动量方向变化率与其垂直，大小
\[
N_{\rm gyro}=\Omega L=I\omega\Omega
=\frac25mR^2\omega\Omega.
\]
轴承相距 $\ell$，
需要一对大小相等、方向相反的横向支承力
形成这一力偶；每端的陀螺力大小为
\[
\boxed{F=\frac{N_{\rm gyro}}{\ell}
=\frac{2mR^2\omega\Omega}{5\ell}
\simeq3.00\times10^2\,\mathrm N
=0.30\,\mathrm{kN}}.
\]
\excheck 两端轴承反力方向相反，
形成一个力偶，具体上、下方向取决于
自转及支架转动的方向，可由
$\bm N=\bm\Omega\times\bm L$ 判定。

\exproblem{287}{46}
陀螺圆盘质量为 $m=15\,\mathrm{kg}$、
半径 $r=5.0\,\mathrm{cm}$，
自转角速度为 $\omega=330\,\mathrm{rad/s}$。
轴承间距为 $\ell=15\,\mathrm{cm}$。
自转轴被迫绕一水平轴作小幅周期摆动，
摆动周期 $T=1.0\,\mathrm s$、
角振幅 $\varphi_m=20^\circ$。
求轴承所受陀螺力的最大值。
\exsolution
圆盘轴向转动惯量
$I=\frac12mr^2$，
自转主角动量 $L=I\omega$。
设自转轴的摆动角
$\varphi(t)=\varphi_m\sin(2\pi t/T)$，
那么轴线方向改变的瞬时角速度
$|\dot\varphi|$ 最大值为 $2\pi\varphi_m/T$。
维持这种方向变化的陀螺力矩最大值
\[
N_{\max}=L\frac{2\pi\varphi_m}{T}
=\frac{\pi mr^2\omega\varphi_m}{T}.
\]
由轴承距 $\ell$，每端最大附加力
\[
\boxed{F_{\max}
=\frac{\pi mr^2\omega\varphi_m}{\ell T}
\simeq90\,\mathrm N=0.090\,\mathrm{kN}},
\]
其中 $\varphi_m=20^\circ$ 须换算为
$0.349\,\mathrm{rad}$。
\excheck 本题的最大陀螺力发生在角位移过零且
角速度模最大的时刻，而不是角位移达到 $\pm\varphi_m$ 时。

\exproblem{288}{46}
轮船以 $v=36\,\mathrm{km/h}$ 沿水平圆弧航行，
轨道半径为 $R=200\,\mathrm m$。
船上沿船长方向布置的飞轮轴
转动惯量为 $I=3.8\times10^3\,\mathrm{kg\,m^2}$，
飞轮转速 $n=300$ 转/分。
求飞轮轴承承受的陀螺力矩大小。
\exsolution
飞轮自转角速度为
$\omega=2\pi n/60$（弧度每秒）。
船体沿水平圆弧改变航向的角速度为
$\Omega=v/R$，其中
$v=36/3.6=10.0\,\mathrm{m/s}$。
飞轮主自转角动量
$L=I\omega$ 沿船长方向，
与航向改变的竖直角速度 $\bm\Omega$ 垂直。
轴承陀螺力矩大小
\[
\boxed{N=I\omega\Omega
=\frac{2\pi nIv}{60R}
=\frac{2\pi(300)(3.8\times10^3)(10)}{60(200)}
\simeq5.97\times10^3\,\mathrm{N\,m}
\approx6.0\,\mathrm{kN\,m}}.
\]
\excheck 力矩的方向取决于转弯方向和飞轮的自转方向，
可用 $\bm N=\bm\Omega\times\bm L$ 确定。

\exproblem{289}{46}
一台机车沿半径 $R=250\,\mathrm m$ 的水平曲线
以速率 $v=50\,\mathrm{km/h}$ 行驶。
汽轮机转轴平行于车轮轴，其自转方向与车轮相同。
汽轮机转子的轴向惯量
$I=240\,\mathrm{kg\,m^2}$，
转速 $n=1500$ 转/分。两轨轨距
$\ell=1.5\,\mathrm m$。
求陀螺效应使铁轨垂直载荷产生的附加变化量。
\exsolution
机车绕竖直方向转弯角速度
$\Omega=v/R$；
汽轮机转子角速度
$\omega=2\pi n/60$。
转轴水平且与航向改变的竖直轴垂直，
故陀螺力矩大小
$N=I\omega\Omega$。
这一力矩由左右轨道上大小相等、方向相反
的\emph{附加竖直反力}形成，
以轨距 $\ell$ 为力偶臂：
\[
\boxed{\Delta F=\frac{N}{\ell}
=\frac{2\pi nIv}{60R\ell}}.
\]
代入
$v=50/3.6\simeq13.89\,\mathrm{m/s}$：
\[
\boxed{\Delta F\simeq1.40\times10^3\,\mathrm N
=1.4\,\mathrm{kN}}.
\]
\excheck 由于汽轮机自转方向与车轮相同，
按原书图示运动方向，\emph{外轨}载荷增加
$\Delta F$、内轨载荷减少同量；
不是两轨总载荷同时增加。
